% Volume VI: Machine Learning Mathematics
% Continuity note for Chapter 32.
% Previous dependencies: Chapter 5's least-squares geometry; Chapter 7's matrix derivatives and shape discipline; Chapters 14-17's Gaussian, Bernoulli, Poisson, exponential-family likelihoods, MLE, and risk; Chapter 31's learning problem, ERM, likelihood-loss links, and generalization.
% New durable concepts: design matrices, linear regression, residuals, normal equations, coefficient interpretation, logistic regression, log odds, sigmoid and softmax classifiers, cross-entropy gradients, generalized linear models, link functions, canonical exponential-family structure, Poisson neural GLMs, decision boundaries, geometric margins, hinge loss, support vectors, and soft-margin SVMs.
% Forward hooks: Chapter 33 regularizes these models and kernelizes them; Chapter 34 turns GLM likelihoods into graphical and latent-variable components; Chapters 46 and 51 reuse Poisson GLMs for spike trains and neural data analysis; deep learning chapters reuse affine scores, logits, softmax, cross-entropy, and margin intuition.
% Notation commitments: Data matrices are X in \mathbb R^{n\times d}, rows are x_i^\top, responses are y_i, coefficients are beta or w, intercepts are b, residuals are r=X\beta-y, sigmoid is \sigma(t)=1/(1+e^{-t}), binary labels use y_i in {0,1} for logistic likelihood and y_i in {-1,+1} for margins, GLM mean is \mu_i=\mathbb E[Y_i\mid x_i], linear predictor is \eta_i=x_i^\top\beta, and link is g(\mu_i)=\eta_i.
% Figure plan: no missing images are included. Proposed ImageGen prompts are recorded in imagegen-figures/ch32-figure-metadata.md.
\section{Chapter 32: Linear Models, GLMs, and Classification}
\label{ch:linear-models-glms-classification}

Chapter 31 defined learning as choosing a function with small expected loss. This chapter studies the simplest function classes that remain genuinely useful: linear predictors, logistic classifiers, generalized linear models, and margin-based classifiers. These models are not just historical warm-ups. They are the local building blocks of deep networks, the standard baselines for neural data analysis, and the clearest place to see how geometry, likelihood, and optimization meet.

The central problem is to understand what a linear score can and cannot express, how different observation models turn that score into a mean or probability, and why classification introduces decision boundaries and margins rather than ordinary residuals alone.

\paragraph{Chapter problem.}
How do linear predictors become regression models, probabilistic classifiers, generalized linear models, and large-margin decision rules?

\paragraph{Section map.}
Section 32.1 revisits linear regression through design matrices, residuals, projection, and Gaussian likelihood. Section 32.2 derives logistic regression and classification boundaries from Bernoulli likelihood. Section 32.3 generalizes the pattern through GLMs: linear predictor, link function, and exponential-family observation model. Section 32.4 introduces margins, hinge loss, support vectors, and support vector machines.

\subsection{32.1 Linear regression revisited}

\paragraph{Guiding question.}
When we fit a linear function to real-valued data, what exactly is being minimized, and what geometry explains the solution?

\paragraph{Beginner-facing intuition.}
Linear regression predicts a real-valued output by adding weighted input features. The word "linear" does not mean the raw data have to lie on a literal line in a plot. It means the prediction is linear in the chosen features. If the features include nonlinear transformations such as \(x^2\), sinusoids, stimulus filters, or neural population summaries, the model is still linear in its coefficients.

The residuals are the errors left after prediction. Least squares chooses coefficients so that the residual vector is as short as possible. Geometrically, it projects the response vector onto the column space of the design matrix.

\paragraph{Formal definitions.}
Let
\[
X=
\begin{bmatrix}
x_1^\top\\
\vdots\\
x_n^\top
\end{bmatrix}
\in\mathbb R^{n\times d},
\qquad
y=
\begin{bmatrix}
y_1\\
\vdots\\
y_n
\end{bmatrix}
\in\mathbb R^n.
\]
The row \(x_i^\top\) contains the \(d\) features for example \(i\). A linear regression model predicts
\[
\widehat y_i=x_i^\top\beta,
\qquad
\beta\in\mathbb R^d.
\]
If an intercept is needed, include a column of ones in \(X\).

The residual vector is
\[
r(\beta)=X\beta-y.
\]
The least-squares objective is
\[
\boxed{
J(\beta)=\frac{1}{2}\|X\beta-y\|_2^2
=\frac{1}{2}\sum_{i=1}^n (x_i^\top\beta-y_i)^2.
}
\]
The factor \(1/2\) is algebraic convenience; it does not change the minimizer.

\paragraph{Core derivation: normal equations and projection.}
Using the matrix derivative convention from Chapter 7,
\[
dJ
=\langle X\beta-y, X\,d\beta\rangle
=\langle X^\top(X\beta-y),d\beta\rangle.
\]
Therefore
\[
\nabla_\beta J=X^\top(X\beta-y).
\]
At an interior optimum,
\[
\boxed{
X^\top X\,\widehat\beta=X^\top y.
}
\]
These are the normal equations.

The geometric meaning is orthogonality of the residual:
\[
X^\top(X\widehat\beta-y)=0.
\]
Thus \(r(\widehat\beta)\) is orthogonal to every column of \(X\). Equivalently, \(X\widehat\beta\) is the orthogonal projection of \(y\) onto the column space \(\operatorname{col}(X)\). If \(X^\top X\) is invertible,
\[
\widehat\beta=(X^\top X)^{-1}X^\top y.
\]
If \(X^\top X\) is singular, solutions may be nonunique, and regularization or pseudoinverses become important in Chapter 33.

\paragraph{Concrete example.}
Fit the one-parameter model \(f_a(x)=a x\) to
\[
(x_i,y_i)=(1,1),(2,3),(3,4).
\]
Here \(X=(1,2,3)^\top\) and \(y=(1,3,4)^\top\). The normal equation is
\[
(1^2+2^2+3^2)\widehat a=1\cdot1+2\cdot3+3\cdot4.
\]
Thus
\[
14\widehat a=19,\qquad \widehat a=\frac{19}{14}\approx1.36.
\]
The fitted predictions are approximately \(1.36,2.71,4.07\). The residuals are approximately \(0.36,-0.29,0.07\), and their weighted orthogonality condition is \(1r_1+2r_2+3r_3=0\).

\paragraph{Computational neuroscience example.}
In a linear receptive-field model, \(x_i\in\mathbb R^d\) may be a vector of stimulus intensities over recent time bins, and \(y_i\) may be a measured membrane voltage or firing-rate estimate. The coefficient vector \(\beta\) is a temporal filter. Least squares fits the filter whose convolution with the stimulus best predicts the response in squared-error average. The coefficient pattern can suggest stimulus sensitivity, but by itself it is not proof of a causal mechanism unless the experimental design supports that interpretation.

\paragraph{AI and machine-learning example.}
Linear regression is a baseline for predicting real-valued targets such as demand, reward, value functions, or embedding coordinates. In deep learning, the final layer before a scalar output is often a linear map from learned features to a prediction. The earlier network learns features; the last layer performs a regression-like affine readout.

\paragraph{Common misconceptions and failure modes.}
First, a large coefficient does not automatically mean a feature is scientifically important; units, scaling, collinearity, and confounding matter. Second, least squares assumes squared-error loss is appropriate; heavy-tailed errors or outliers can dominate the fit. Third, invertibility of \(X^\top X\) is not guaranteed, especially when features are redundant or \(d>n\). Fourth, a good fit on training data does not establish causal effects.

\paragraph{Exercises.}
\begin{enumerate}[label=\textbf{32.1.\arabic*.}, leftmargin=*]
\item \textbf{Mechanical.} For \(X\in\mathbb R^{100\times 5}\), \(y\in\mathbb R^{100}\), and \(\beta\in\mathbb R^5\), state the shapes of \(X\beta\), \(X^\top y\), and \(X^\top X\).
\item \textbf{Proof or derivation.} Starting from \(J(\beta)=\frac12\|X\beta-y\|_2^2\), derive \(\nabla_\beta J=X^\top(X\beta-y)\).
\item \textbf{Numerical/computational.} Fit \(f_a(x)=a x\) to \((1,2),(2,2),(3,5)\) by solving the one-parameter normal equation.
\item \textbf{Interpretation.} In a neural receptive-field regression, what does it mean for the residual vector to be orthogonal to the stimulus-feature columns?
\item \textbf{Misconception/debugging.} A student says, "The feature with the largest coefficient causes the response most strongly." Give two reasons this interpretation can fail.
\end{enumerate}

\subsection{32.2 Logistic regression and classification}

\paragraph{Guiding question.}
How can a linear score be turned into a probability of a binary class, and how does the loss follow from likelihood?

\paragraph{Beginner-facing intuition.}
Classification is not ordinary regression with labels renamed as numbers. A binary label is constrained to \(0\) or \(1\), while a linear score \(w^\top x+b\) can be any real number. Logistic regression solves this mismatch by using the linear score as a log-odds, then mapping it through the sigmoid function to obtain a probability.

The decision boundary remains linear: points with \(w^\top x+b=0\) have predicted probability \(1/2\). The sigmoid changes how confident the model is away from that boundary.

\paragraph{Formal definitions.}
For binary labels \(Y\in\{0,1\}\), define the score
\[
s(x)=w^\top x+b,
\qquad w\in\mathbb R^d,\ b\in\mathbb R.
\]
The sigmoid function is
\[
\sigma(t)=\frac{1}{1+e^{-t}}.
\]
Logistic regression models
\[
\boxed{
P(Y=1\mid X=x)=\sigma(w^\top x+b).
}
\]
The log-odds are linear:
\[
\log\frac{P(Y=1\mid X=x)}{P(Y=0\mid X=x)}
=w^\top x+b.
\]
Given data \((x_i,y_i)_{i=1}^n\), the Bernoulli negative log-likelihood is
\[
\boxed{
L(w,b)=
-\sum_{i=1}^n
\left[
y_i\log p_i+(1-y_i)\log(1-p_i)
\right],
\qquad
p_i=\sigma(w^\top x_i+b).
}
\]
This is binary cross-entropy.

\paragraph{Core derivation: gradient of logistic loss.}
Let \(s_i=w^\top x_i+b\), \(p_i=\sigma(s_i)\), and
\[
\ell_i=-y_i\log p_i-(1-y_i)\log(1-p_i).
\]
The sigmoid derivative is
\[
\sigma'(s)=\sigma(s)(1-\sigma(s)).
\]
Differentiate \(\ell_i\) with respect to \(s_i\):
\[
\frac{d\ell_i}{ds_i}
=-\frac{y_i}{p_i}p_i(1-p_i)
+\frac{1-y_i}{1-p_i}p_i(1-p_i)
=p_i-y_i.
\]
Since \(ds_i/dw=x_i\) and \(ds_i/db=1\),
\[
\boxed{
\nabla_w L=\sum_{i=1}^n (p_i-y_i)x_i,
\qquad
\frac{\partial L}{\partial b}=\sum_{i=1}^n(p_i-y_i).
}
\]
The gradient is prediction minus label, multiplied by the input. This pattern reappears in softmax classifiers and deep-network backpropagation.

\paragraph{Concrete example.}
Let \(w=(1,-1)\), \(b=0\), and \(x=(2,1)\). Then
\[
s=w^\top x=1,
\qquad
P(Y=1\mid x)=\sigma(1)\approx0.731.
\]
For \(x=(1,2)\), the score is \(-1\), so the predicted probability is \(\sigma(-1)\approx0.269\). The decision boundary is
\[
x_1-x_2=0,
\]
or \(x_1=x_2\). Points with \(x_1>x_2\) are classified as class \(1\) at threshold \(1/2\).

\paragraph{Computational neuroscience example.}
A logistic model can predict whether a subject chose "right" or "left" from neural activity:
\[
P(Y=\text{right}\mid x)=\sigma(w^\top x+b),
\]
where \(x\) is a vector of spike counts or latent neural factors. The weight \(w_j\) measures how neuron or feature \(j\) changes the log-odds after accounting for other features. It should not be read as a causal effect unless the features were experimentally manipulated or confounding was controlled.

\paragraph{AI and machine-learning example.}
For \(K\)-class classification, softmax regression uses scores
\[
s_k(x)=w_k^\top x+b_k
\]
and probabilities
\[
p_k(x)=\frac{e^{s_k(x)}}{\sum_{\ell=1}^K e^{s_\ell(x)}}.
\]
The cross-entropy loss for label \(y\) is \(-\log p_y(x)\). Deep networks usually replace \(x\) by learned features \(h_\theta(x)\), but the final classification layer is still a linear score plus softmax.

\paragraph{Common misconceptions and failure modes.}
First, logistic regression is linear in log-odds, not linear in predicted probability. Second, a probability near \(0.9\) is a model output, not a guaranteed frequency unless the model is calibrated. Third, a linear decision boundary in the chosen features can be nonlinear in raw inputs if the features are nonlinear. Fourth, class imbalance can make accuracy misleading; likelihood, calibration, precision-recall, and task-specific losses may matter.

\paragraph{Exercises.}
\begin{enumerate}[label=\textbf{32.2.\arabic*.}, leftmargin=*]
\item \textbf{Mechanical.} Compute \(\sigma(0)\), \(\sigma(2)\), and \(\sigma(-2)\) to three decimal places.
\item \textbf{Proof or derivation.} Derive \(d\ell_i/ds_i=p_i-y_i\) for binary cross-entropy and sigmoid output.
\item \textbf{Numerical/computational.} For \(w=(2,1)\), \(b=-1\), compute the score, predicted probability, and threshold-\(1/2\) class for \(x=(1,0)\) and \(x=(0,2)\).
\item \textbf{Interpretation.} If a neural-choice model has \(w_j>0\) for neuron \(j\), what does that mean in log-odds terms, holding other features fixed?
\item \textbf{Misconception/debugging.} A classifier has \(95\%\) accuracy on a dataset where \(95\%\) of examples are class \(0\). Explain why this may be a poor classifier.
\end{enumerate}

\subsection{32.3 Generalized linear models}

\paragraph{Guiding question.}
How can one linear predictor support Gaussian responses, binary labels, counts, and other observation types?

\paragraph{Beginner-facing intuition.}
A generalized linear model, or GLM, has three pieces. First, compute a linear predictor \(\eta=x^\top\beta\). Second, pass it through a link or inverse-link function to obtain a mean \(\mu\). Third, choose a noise distribution appropriate for the observation type. The model is linear in the predictor, but the mean response can be nonlinear in the inputs.

This is exactly the pattern needed for neural data. Voltages are continuous, choices are binary, spike counts are nonnegative integers, and reaction times are positive. A single Gaussian least-squares model cannot respect all those supports.

\paragraph{Formal definitions.}
For observations \(Y_i\) with features \(x_i\in\mathbb R^d\), a GLM specifies:
\[
\eta_i=x_i^\top\beta,
\]
\[
\mu_i=\mathbb E[Y_i\mid x_i],
\]
and a link function \(g\) such that
\[
\boxed{
g(\mu_i)=\eta_i.
}
\]
Equivalently, \(\mu_i=g^{-1}(\eta_i)\). The conditional distribution of \(Y_i\mid x_i\) is chosen from an exponential family:
\[
p(y\mid \eta)
=h(y)\exp\bigl(\eta T(y)-A(\eta)\bigr)
\]
in a scalar canonical form, or a vector generalization.

\begin{center}
\small
\setlength{\tabcolsep}{3pt}
\renewcommand{\arraystretch}{1.2}
\begin{tabular}{p{0.16\linewidth}p{0.18\linewidth}p{0.18\linewidth}p{0.22\linewidth}p{0.18\linewidth}}
\toprule
Observation & Distribution & Mean \(\mu\) & Link \(g(\mu)\) & Typical use \\
\midrule
real value & Gaussian & \(\mu\in\mathbb R\) & identity \(\mu\) & regression, voltage \\
binary label & Bernoulli & \(0<\mu<1\) & logit \(\log(\mu/(1-\mu))\) & choices, classification \\
count & Poisson & \(\mu>0\) & log \(\log\mu\) & spike counts, events \\
positive value & Gamma & \(\mu>0\) & inverse or log & waiting times, rates \\
\bottomrule
\end{tabular}
\end{center}

\paragraph{Core derivation: Poisson GLM likelihood and gradient.}
For count data, use
\[
Y_i\mid x_i\sim\operatorname{Pois}(\lambda_i),
\qquad
\lambda_i=\exp(x_i^\top\beta).
\]
The log link \(\log\lambda_i=x_i^\top\beta\) ensures \(\lambda_i>0\). The log-likelihood is
\[
\ell(\beta)
=\sum_{i=1}^n
\left[
y_i\log\lambda_i-\lambda_i-\log(y_i!)
\right].
\]
Substituting \(\log\lambda_i=x_i^\top\beta\) gives
\[
\ell(\beta)
=\sum_{i=1}^n
\left[
y_i x_i^\top\beta-\exp(x_i^\top\beta)-\log(y_i!)
\right].
\]
Differentiating,
\[
\boxed{
\nabla_\beta \ell(\beta)
=\sum_{i=1}^n x_i\left(y_i-\lambda_i\right).
}
\]
The score has the form observed minus predicted count, weighted by features. This mirrors the logistic-gradient pattern \(p_i-y_i\), with signs depending on whether one maximizes likelihood or minimizes negative log-likelihood.

\paragraph{Concrete example.}
Suppose a Poisson GLM has one feature \(x\) and parameters \(\beta_0=\log 2\), \(\beta_1=0.5\):
\[
\lambda(x)=\exp(\beta_0+\beta_1x).
\]
At \(x=0\), \(\lambda(0)=2\). At \(x=2\),
\[
\lambda(2)=\exp(\log 2+1)=2e\approx5.44.
\]
If \(y=4\) spikes are observed at \(x=2\), the log-likelihood contribution is
\[
4\log(5.44)-5.44-\log(4!).
\]
The model respects count support because it assigns probabilities only to nonnegative integers.

\paragraph{Computational neuroscience example.}
A common neural GLM predicts spike count or spike probability from stimulus filters, spike-history filters, and coupling terms:
\[
\lambda_i(t)=\exp\bigl(k_i^\top s(t)+h_i^\top r_i(t-\cdot)+\sum_{j\neq i}c_{ij}^\top r_j(t-\cdot)+b_i\bigr).
\]
Here the linear predictor combines stimulus drive, refractoriness or bursting history, and other neurons' recent spikes. The exponential inverse link converts this predictor into a nonnegative conditional intensity or count rate. The model is statistical: coupling weights describe predictive dependence in the fitted model and require care before being interpreted as synapses.

\paragraph{AI and machine-learning example.}
Logistic regression and softmax regression are GLMs for Bernoulli and categorical outputs. Poisson GLMs model event counts, click counts, word counts in simple text models, and demand arrivals. GLMs also appear inside larger systems as output heads: a neural network can learn features \(h_\theta(x)\), while a GLM-like output layer chooses an observation distribution appropriate to the target.

\paragraph{Common misconceptions and failure modes.}
First, the link function is not the loss. It connects the mean to the linear predictor. Second, GLMs are not limited to linear raw features; nonlinear basis functions can be included in \(x\). Third, the observation distribution matters. Using a Gaussian model for counts can predict impossible negative values. Fourth, overdispersion in neural counts can violate the Poisson mean-equals-variance assumption, requiring richer models.

\paragraph{Exercises.}
\begin{enumerate}[label=\textbf{32.3.\arabic*.}, leftmargin=*]
\item \textbf{Mechanical.} For a Bernoulli GLM with logit link and \(\eta=1.5\), compute the mean \(\mu\).
\item \textbf{Proof or derivation.} Derive the Poisson GLM gradient \(\nabla_\beta \ell(\beta)=\sum_i x_i(y_i-\lambda_i)\).
\item \textbf{Numerical/computational.} A Poisson GLM has \(\lambda(x)=\exp(0.2+0.3x)\). Compute \(\lambda(0)\), \(\lambda(2)\), and the log-likelihood contribution for \(y=3\) at \(x=2\).
\item \textbf{Interpretation.} In a spike-count GLM, why is the log link useful?
\item \textbf{Misconception/debugging.} A student says, "A GLM is just linear regression with a fancier name." Explain the roles of the response distribution and link function.
\end{enumerate}

\subsection{32.4 Margins and support vector machines}

\paragraph{Guiding question.}
For classification, why might we prefer a boundary that separates classes with a large margin rather than merely one that classifies the training data correctly?

\paragraph{Beginner-facing intuition.}
Many different lines or hyperplanes can separate a training set. A margin measures how far the closest correctly classified points are from the boundary. A large margin is a geometric form of confidence: small perturbations of the input are less likely to flip the decision. Support vector machines choose a separating hyperplane by maximizing this margin, or by trading margin against classification violations when data are not perfectly separable.

\paragraph{Formal definitions.}
Use labels \(y_i\in\{-1,+1\}\). A linear classifier has score
\[
s(x)=w^\top x+b
\]
and predicts \(\operatorname{sign}(s(x))\). A point is correctly classified when
\[
y_i(w^\top x_i+b)>0.
\]
The geometric margin of point \(i\) is
\[
\gamma_i=\frac{y_i(w^\top x_i+b)}{\|w\|_2}.
\]
This is the signed distance from \(x_i\) to the hyperplane \(w^\top x+b=0\). The margin of the classifier on the sample is \(\min_i\gamma_i\).

For separable data, one can scale \(w,b\) so that the closest points satisfy
\[
y_i(w^\top x_i+b)=1.
\]
The hard-margin SVM solves
\[
\boxed{
\min_{w,b}\frac12\|w\|_2^2
\quad\text{subject to}\quad
y_i(w^\top x_i+b)\geq1\ \text{for all }i.
}
\]
Minimizing \(\|w\|_2\) under these constraints maximizes the geometric margin.

\paragraph{Core derivation: why the margin width is \(2/\|w\|\).}
The two canonical margin planes are
\[
w^\top x+b=1
\qquad\text{and}\qquad
w^\top x+b=-1.
\]
Choose points \(x_+\) and \(x_-\) on these two planes that differ only in the normal direction. The normal unit vector is \(w/\|w\|\). The distance \(d\) between the planes must satisfy
\[
x_+-x_-=d\frac{w}{\|w\|}.
\]
Subtract the plane equations:
\[
w^\top(x_+-x_-)=2.
\]
Substitute the normal displacement:
\[
w^\top\left(d\frac{w}{\|w\|}\right)=d\|w\|=2.
\]
Thus
\[
\boxed{d=\frac{2}{\|w\|}.}
\]
Maximizing the distance between the margin planes is therefore equivalent to minimizing \(\|w\|\), after the scale convention \(y_i(w^\top x_i+b)\geq1\) is imposed.

\paragraph{Soft margins and hinge loss.}
Real data are often not linearly separable. The hinge loss
\[
\ell_{\mathrm{hinge}}(w,b;x_i,y_i)
=\max\{0,1-y_i(w^\top x_i+b)\}
\]
penalizes points inside the margin or on the wrong side of the boundary. A soft-margin objective is
\[
\boxed{
\min_{w,b}
\frac12\|w\|_2^2
+C\sum_{i=1}^n
\max\{0,1-y_i(w^\top x_i+b)\}.
}
\]
The parameter \(C>0\) controls the tradeoff between a large margin and training violations. Points with nonzero hinge loss, or exactly on the margin in the constrained form, are the support vectors; they determine the fitted boundary.

\paragraph{Concrete example.}
Consider two points on the real line:
\[
x_1=-1,\ y_1=-1,\qquad x_2=1,\ y_2=+1.
\]
The classifier \(s(x)=x\) has \(w=1,b=0\). Both constraints hold:
\[
y_1s(x_1)=(-1)(-1)=1,\qquad y_2s(x_2)=1.
\]
The boundary is \(x=0\), the margin planes are \(x=-1\) and \(x=1\), and the margin width is \(2/\|w\|=2\). Both training points are support vectors.

\paragraph{Computational neuroscience example.}
A linear SVM can classify stimulus category from population activity vectors. The normal vector \(w\) defines a decoding direction in neural state space. A large margin means the two classes are separated with a buffer in that representation. This is useful for decoding, but the margin direction should not automatically be interpreted as the axis used by the brain; it is the axis chosen by the decoder and data preprocessing.

\paragraph{AI and machine-learning example.}
Before deep learning became dominant, SVMs with linear or kernel features were standard high-performance classifiers for text, vision features, and bioinformatics. Margin ideas remain important in deep learning: logits, contrastive losses, metric learning, adversarial robustness, and representation geometry all use variants of "make correct scores separated from incorrect scores."

\paragraph{Common misconceptions and failure modes.}
First, scaling \(w,b\) changes raw scores but not the decision boundary; the SVM constraints fix scale so margin comparisons are meaningful. Second, an SVM output is not automatically a calibrated probability. Third, not every training point determines the boundary; support vectors do. Fourth, a large training margin can still fail under distribution shift or a poor feature representation.

\paragraph{Exercises.}
\begin{enumerate}[label=\textbf{32.4.\arabic*.}, leftmargin=*]
\item \textbf{Mechanical.} For \(w=(3,4)\), \(b=-1\), \(x=(1,1)\), and \(y=+1\), compute the score and geometric margin.
\item \textbf{Proof or derivation.} Derive the distance from a point \(x_0\) to the hyperplane \(w^\top x+b=0\) as \(|w^\top x_0+b|/\|w\|_2\).
\item \textbf{Numerical/computational.} For \(y_i(w^\top x_i+b)\) values \(2.0,1.0,0.4,-0.5\), compute the hinge losses.
\item \textbf{Interpretation.} In a neural population decoder, what does a small margin suggest about class separability and sensitivity to noise?
\item \textbf{Misconception/debugging.} A student says, "If a classifier separates the training data, all separating hyperplanes are equally good." Explain the margin-based objection.
\end{enumerate}

\paragraph{Closing bridge.}
Linear regression, logistic regression, GLMs, and SVMs show how far one can go with linear scores plus appropriate losses and observation models. Chapter 33 adds explicit control of complexity through regularization and then replaces explicit finite feature vectors by kernels. That path leads naturally to Gaussian processes, where the unknown function itself is assigned a probability distribution.

\clearpage
